\subsection{Finite controls, random stopping and physical resolution}
The next results add two resource converses to the preceding constructive
inverse. The geometric regularity remains $s=6+\beta$, and all uniform
finite global decisions still require the bounded periodic class and
known positive patch margin. No new geometric response is supplied.

\begin{theorem}[Attempted bits and calibration resolution]
\label{thm:resources}
The adaptive reciprocal experiment has a realization with dyadic centers
and scales, using $O(\log(C/\nu)+\log(C/\delta))$ binary digits per
batch command and the same attempted-bit bound
\eqref{eq:attempts}. On a fixed nondegenerate physical subclass, any
procedure whose only table-dependent data are collision bits and whose
stopping decision uses only its past bits and independent randomness
has worst-case expected attempt count at least
$c\nu^{-1/(s-2)}$ at uniform confidence $3/4$.

Moreover, under table-, command- and draw-dependent bounded-position-error
implementations,
uniform centered $C^2$ accuracy $\nu$ with confidence greater than
$1/2$ for short commands requires
\[
                            r\le C\nu^{s/(s-2)}.
\]
Here $r$ is the allowed position error, durations and directions may
be exact, and command lengths have any fixed bound less than the
component separation. Above a constant multiple of this scale,
two distinct physical tables admit identical complete adaptive bit
experiments under two possibly different admissible implementations.
No identical common-noise device is asserted. More attempted
bits do not remove this obstruction. The sufficient combined tolerance
$\ell+\tau+t\alpha\le c\nu^{s/(s-2)}$ remains as before.
\end{theorem}

Theorem~\ref{thm:digital} proves the finite-description claim and gives
a separate conservative numerical cost. Theorem~\ref{thm:expected}
proves the expected-stopping claim by a prefix-free transcript argument,
with the centering gauge explicit. Theorem~\ref{thm:calibration-floor}
proves the physical-resolution converse through a pointwise common-response
coupling of nearby convex swept bodies. Neither converse prices
apparatus production, and neither claims sharp logarithmic factors.
The deterministic-cap lower bound is retained independently: it is
not used as an unproved substitute for the sequential result.
